\input{preamble}

\begin{document}

\title{Reading the Masters}
\maketitle

\phantomsection
\label{section-phantom}

\noindent
Personal notes on important mathematical
articles I read, one paragraph at a time.

\medskip
\tableofcontents

\section{Saint-Donat: first impressions}
\label{section-saint-donat}

\noindent
B. Saint-Donat, \textit{Projective Models
of K-3 Surfaces} \cite{saint-donat}.
Notes from the introduction, 9 September
2026. Here the K3 surfaces are projective
over $\mathbb{C}$.

\medskip\noindent
\textit{My first impressions, with minimal
corrections.}
I gather that the cool thing about
Saint-Donat's paper is that a line bundle
$L$ with $L^2>0$ provides a morphism to
projective space if $|L|$ is nonempty and
has no fixed components. I first understood
this as a projective embedding, but that
requires more. If I remember correctly,
Nakai--Moishezon says that $L$ is ample if
$L^2>0$ and $L\cdot C>0$ for every
irreducible curve $C$. This condition on
curves still matters on a K3; ampleness is
also different from very ampleness.

Saint-Donat also says he derives some other
theorem involving no s\'e qu\'e
non-hyperelliptic curves. I do not yet
understand that part.

And there was a third thing\ldots

\medskip\noindent
\textbf{AI clarification.}
The second point is an analogue of the
Enriques--Petri theorem about equations of
non-hyperelliptic canonical curves. For
non-hyperelliptic K3 projective models,
Saint-Donat proves
generation of the homogeneous ideal by
quadrics outside explicitly described
exceptions; see
\cite[Theorem 7.2]{saint-donat}.
The third point is that, if $L$ is ample,
then $L^{\otimes n}$ is basepoint-free for
$n\geq 2$ and very ample for $n\geq 3$;
see \cite[Theorem 8.3]{saint-donat}.

\medskip\noindent
\textbf{My question.}
So every line bundle on every variety
defines a birational map to projective
space?

\medskip\noindent
\textbf{AI answer.}
No. Let $X$ be an integral projective
variety over $\mathbb{C}$ and $L$ a line
bundle. If $H^0(X,L)=0$, the construction
using global sections gives no map.
Otherwise, a basis $s_0,\ldots,s_N$ of
$H^0(X,L)$ gives a rational map
$$
\begin{aligned}
\varphi_{|L|}:X
&\dashrightarrow\mathbb{P}^N,\\
x&\longmapsto[s_0(x):\cdots:s_N(x)].
\end{aligned}
$$
This formula is defined wherever the
sections do not all vanish. If $L$ is
basepoint-free, it defines a morphism on
all of $X$. If $L$ is very ample, the
morphism is a closed embedding.

Being \textit{birational onto its image}
means that the map restricts to an
isomorphism between dense open subsets
of $X$ and its image. This is an extra
property, not part of the meaning of
``rational map.''

For example, the trivial line bundle
$\mathcal{O}_{\mathbb{P}^1}$ has only
constant global sections. Its complete
linear system maps $\mathbb{P}^1$ to a
point, so the map is a morphism but is
not birational onto its image.

\section{Apéry: Fano threefolds and links}
\label{section-apery-fano-links}

\noindent
Roger Apéry, \textit{Étude de certaines
hypersurfaces}, Chapter III of his 1947
thesis. In modern language, Apéry studies
Picard-rank-one Fano threefolds and
birational maps obtained by blowing up a
curve on each side. The striking point is
that the numerical relations he derives
are essentially the numerical theory of
type-II Sarkisov links: his ``curves of the
second kind'' are the flopping curves.
Compare Section~\ref{section-mmp-and-sarkisov},
especially Theorem~\ref{theorem-sarkisov}.
There the HP states the modern Sarkisov
program: birational maps of Mori fibre
spaces are decomposed into Sarkisov links.
Apéry is studying a concrete rank-one Fano
case of this picture, where one blows up a
curve, passes through the flopping curves,
and contracts a curve on the other side.
His tables also contain the genus-$12$ Fano
threefold $V_{22}$ and locate the step in
Fano's argument that missed it.

\medskip\noindent
\textbf{Sergey's reconstruction.}
Sergey Galkin's September 2026 working
draft makes Chapter III readable in modern
terms: it gives the French text beside an
English translation, supplies a dictionary
from Apéry's terminology to current Fano
geometry, rewrites the birational
constructions as blow-up--flop--blow-down
Sarkisov links, and checks Apéry's fifty
table rows against the modern
classification. All fifty rows are
accounted for by the modern
classification, although Apéry's tables
are incomplete; Sergey also identifies the
few misprints and one false non-existence
statement in the chapter.

\bibliography{my}
\bibliographystyle{amsalpha}

\end{document}
